% English-first editorial revision; the frozen Russian source is not synchronized.
This article examined the transition from the duration of an isolated
AI-assisted task to the makespan of a fixed set of interdependent software
tasks. For the human--agent cell studied here, the answer to the research
question is conditional: a configuration reduces makespan relative to
sequential no-AI work only when a feasible schedule completes the same workload
to the same acceptance criterion faster than the baseline. Achievable speedup
depends jointly on task- and mode-specific durations, dependencies and phase
profiles, the number of agent streams and the rules governing their use,
integration overhead, and the sequential resource of human attention. It is
therefore a property of a fully specified process, not an isolated
characteristic of the AI tool or the nominal agent count.

The M0--M4 hierarchy makes this conclusion progressively less idealized. It
moves from divisible aggregate work to indivisible tasks, precedence
constraints, configuration-dependent overhead, and an exact phase schedule in
which autonomous-agent work alternates with requests for one developer's
attention. At M0, the exact aggregate speedup condition is
$P>\bar{k}_w$; under the stated identical-stream and zero-overhead assumptions,
the M1--M2 list-scheduling bounds give sufficient speedup guarantees accounting
for task indivisibility and precedence. At M4, the structural lower bound $B_4$
combines agent-stream load, the original critical path, and total human work.
The bound is a necessary diagnostic, but it need not be attainable: repeated
human service and local retention of a waiting task's agent slot can create
mixed resource chains that are absent from its three branches. The
maximum-weight paths in feasible resource-augmented phase graphs capture those
chains, and minimizing their length gives the optimal makespan exactly.
Sequential human work also imposes a speedup ceiling; after it becomes active,
additional agents cannot substitute for shorter specification, review,
correction, or acceptance phases.

The reproducible EXP-00--EXP-08 package computationally checks these
distinctions within the formal model rather than calibrating a particular tool
or team. Most importantly, in the frozen 90-scenario EXP-01 grid, a strict
certified gap $T^*>B_4$ occurred in 30 cases, so the lower bound cannot be used
as a universal estimate of achievable makespan. The sensitivity experiments
also separate available capacity from configured parallelism: an unused
additional stream cannot worsen the optimum when phase durations and the
overhead multipliers $C$ and $\gamma$ are held fixed, whereas a change in
configuration may increase integration or attention overhead and create an
interior optimum over $P$. Likewise, finer decomposition can expose parallelism
while adding specification and integration boundaries, producing nonmonotonic
makespan profiles. Complementarity between parallelism and granularity therefore
does not imply that either more streams or more tasks is unconditionally
beneficial. All reported frequencies describe prespecified synthetic matrices
or generated samples, not the prevalence of these effects in real projects.

In practice, the model should be used to compare complete work-organization
scenarios rather than to select the largest available agent count. Tasks and
operating modes must be placed on a common scale, the workload and acceptance
criterion must be fixed, and the DAG, phase profiles, and overhead assumptions
must be stated. The active branch of $B_4$ identifies the constraint governing
the structural lower bound, but only a resource-feasible schedule establishes
an attainable makespan. Sensitivity analysis is therefore part of the
comparison, not an optional addition. When candidate interventions are encoded
as complete alternative scenarios, sensitivity comparisons can indicate
whether makespan is more responsive to task execution time, critical-path
length, human-attention demand, coordination overhead, or decomposition.

These conclusions remain limited to the assumed architecture: a fixed, known
DAG; one developer; homogeneous agent streams; nonpreemptive phases; and local
slot retention. Quality and computational cost are represented only when they
alter time to an accepted result, and the model does not replace a causal study
of AI effectiveness. Empirical use requires local measurement of agent and
human-attention intervals, waiting, rework, and the dependence of overhead on
configuration. Extensions to dynamic graph revelation, adaptive decomposition,
stochastic durations, and other orchestration architectures remain future
work. Until such validation is available, the contribution is structural and
operationally interpretable: within the model, useful parallelism must be
assessed through the coordinated design of tasks, dependencies, phases,
resources, and acceptance rules for the human--agent process as a whole.
